\begin{abstract}
AI-based inpainting can remove, replace, or insert content in a photograph while leaving much of the scene unchanged. Detecting such edits requires both an image-level decision and a spatial explanation of the decision. Small edits pose a particular challenge: resizing a photograph can erase subtle local evidence, while a close crop may lack the scene context needed to interpret it. We investigate a dual-view approach that combines whole-image context with overlapping native-resolution crops for detection and localization from a single input image. The detector combines pretrained ConvNeXt and frozen DINOv2 representations with residual features extracted by fixed and constrained filters. A cross-attention feature pyramid produces an image score and a pixel-level edit map; tiled inference merges local predictions across the original image. We study both an ensemble of earlier checkpoints and a separate joint whole-image and crop training configuration. On an inspected 100-image development subset covering five editor groups, the ensemble achieves 70.0\% image accuracy, 0.7452 AUROC, and 34.89\% mean localization intersection over union (IoU); IoU falls to 8.94\% for edits occupying less than 2\% of the image. Joint-view training reduces false alarms on genuine images in a separate comparison, but also misses more edited images. A separate 10K training summary reports 91.20\% accuracy and 47.30\% mean IoU; its aggregate values require reconciliation with the recorded counts and evaluation artifacts. These findings motivate source-disjoint testing, matched baselines, and explicit evaluation by edit size and generator before establishing comparative performance.
\end{abstract}
